\subsection*{Value, Value Propositions, and Customer Needs}

\textbf{Value} can be defined as the set of benefits a product or service provides to a 
customer, relative to their needs. Researchers such as Ruth N.\ Bolton argue that perceived 
value depends heavily on context and on the individual evaluating it --- different individuals 
hold different priorities, which introduces inherent ambiguity into any valuation.

\subsubsection*{Value Creation Models}

Porter conceptualised value creation as a linear chain: a supplier produces something and 
passes it downstream to the customer. Wikström and Normann proposed a more dynamic 
alternative --- the \textbf{value star} --- in which a web of interconnected actors 
contribute simultaneously. Uber illustrates this well: the driver, the company, and the 
passenger each create and receive value at the same time.

\subsubsection*{Perceived Value in Services}

Perceived value for a service is shaped by four factors: quality, price, needs, and 
expectations.

\subsubsection*{Perceived Benefits vs.\ Perceived Sacrifice}

Perceived benefits encompass the physical attributes of a product (quality, features), 
service attributes (support, delivery, assistance), technical support, perceived quality 
indicators, and performance in the customer's actual use case.

Perceived sacrifice includes all costs the customer faces, not merely the purchase price. 
These include time spent purchasing or learning to use the product, switching costs, and 
any other friction involved.

Formally:
\[
    \text{Perceived Value} = \frac{\text{Perceived Benefits}}{\text{Perceived Sacrifice}}
\]

Ravald and Grönroos argue that value increases either when benefits increase or when 
sacrifice decreases. A firm can therefore increase perceived value by enhancing what the 
customer receives \textit{or} by reducing the friction of buying. Amazon is an instructive 
example: product quality is broadly equivalent to competitors, yet removing shipping 
friction substantially increases perceived value.

\subsubsection*{Product as a ``Value Carrier''}

A product should always serve as a \textbf{value carrier}. A firm achieves sustainable 
competitive advantage when customers believe its offering delivers greater net perceived 
benefits than any competing alternative.

\subsubsection*{Relationship Value}

Holmlund and Kock introduce the idea that the buyer--seller relationship itself carries 
value, through attributes such as trust, commitment, reliability, long-term cooperation, 
and reduced uncertainty. Customer value should therefore be understood as:
\[
    \text{Total Value} = \text{Episode Value} + \text{Relationship Value}
\]

\textbf{Episode value} refers to the value derived from a specific transaction, determined 
by quality, price, needs, and expectations. \textbf{Relationship value} refers to the 
value that accrues from an ongoing relationship: trust, stability, familiarity, and 
reduced transaction costs. Importantly, a strong relationship can compensate for a 
relatively weak individual transaction.

\subsubsection*{Holmlund \& Kock's Three Quality Dimensions}

Holmlund and Kock identify three dimensions of quality in business relationships:

\begin{itemize}
    \item \textbf{Technical Quality} --- whether the offering meets specifications, fits 
    the buyer's needs, and is accompanied by appropriate documentation.
    \item \textbf{Functional Quality} --- the quality of communication, cooperation, and 
    interpersonal interaction throughout the relationship.
    \item \textbf{Economic Quality} --- all financial costs associated with the 
    relationship, including coordination costs such as meetings and production costs.
\end{itemize}

\subsubsection*{Service Quality Theory}

Grönroos distinguishes between \textbf{technical quality} (what the customer receives) 
and \textbf{functional quality} (how the service is delivered). A failure in technical 
quality can, in some circumstances, be offset by sufficiently strong functional quality.

\subsubsection*{Key Qualities Customers Prioritise}

Research identifies three qualities customers prioritise above others: competitiveness, 
reliability, and adaptability.

\subsubsection*{Economic Quality for Each Party}

For the \textit{buyer}, economic quality means the relationship must be profitable --- 
the value received must exceed the costs incurred. For the \textit{supplier}, economic 
quality means the price paid by the customer must cover the costs of serving them. 
Notably, price is not the primary driver of value; expertise and reliability typically 
matter more.

\subsubsection*{The Diachronic Nature of Value Creation}

Value is not created in a single moment but develops through repeated interactions over 
time. Long-term outcomes matter, and the effects of a relationship accumulate --- this 
is the \textbf{diachronic} nature of value creation.

\subsection*{Customer Value Propositions}

A \textbf{customer value proposition} is a statement that conveys what a brand does and 
how it differs from competitors. Emerging in the 1980s, it is now considered a central 
organising principle, strongly correlated with better business performance across 
marketing strategy, product development, customer relationships, and competitive 
positioning. A value proposition answers three questions:

\begin{itemize}
    \item What benefits do customers receive?
    \item Why is the firm's offering superior to competitors?
    \item How does the company intend to deliver that value?
\end{itemize}

Practically, a value proposition aids in converting potential leads into customers, 
pitching to investors, and communicating clearly what the firm actually does.

\subsubsection*{Jobs to Be Done Theory}

Developed by Harvard Business School Professor Clayton Christensen, \textbf{Jobs to Be 
Done} theory holds that customers do not purchase products based on attributes or 
demographics --- they \textit{hire} products and services to accomplish a specific job. 
The Warby Parker case illustrates this: the founders identified that glasses were 
artificially expensive due to a near-monopoly (the same company controlled LensCrafters, 
Pearle Vision, Ray-Ban, and Oakley), and repositioned their offering around the job of 
providing affordable, socially conscious eyewear.

\subsubsection*{How to Create a Value Proposition}

A value proposition can be structured around four questions:

\begin{enumerate}
    \item \textit{For} --- what is my brand offering?
    \item \textit{I will make a difference by} --- what job does the customer hire my 
    brand to do?
    \item \textit{Among all} --- what companies and products compete with my brand for 
    this job?
    \item \textit{Because} --- what sets my brand apart from those competitors?
\end{enumerate}

\subsection*{Types of Customer Needs}

There are three categories of customer needs:

\begin{enumerate}
    \item \textbf{Functional Needs} --- practical, task-oriented requirements.
    \item \textbf{Social Needs} --- the desire to be perceived positively by others.
    \item \textbf{Emotional Needs} --- feelings and experiences the customer seeks.
\end{enumerate}

\subsubsection{Pain Points vs Desired Gains}\mbox{}
\subsubsection{Problem Awareness Levels}\mbox{}
\footcite{christensen2016know}
